\documentclass[14pt,a4paper]{extarticle}
\usepackage[T2A]{fontenc}
\usepackage[utf8]{inputenc}
\usepackage[english,russian]{babel}
\usepackage{lmodern}
\usepackage{geometry}
\usepackage{array}
\usepackage{tabularx}
\usepackage{longtable}
\usepackage{ltablex}
\usepackage{booktabs}
\usepackage{hyperref}
\usepackage{indentfirst}
\usepackage{microtype}
\usepackage{graphicx}
\usepackage{setspace}
\usepackage{listings}
\usepackage{graphicx}
\usepackage{caption}
\usepackage{float}
\usepackage{tocloft}
\usepackage{titlesec}
\usepackage{newtxtext,newtxmath}

\titleformat{\section}        % форматируем \section
  {\normalfont\bfseries}      % стиль: обычный + жирный
  {\thesection.}              % номер секции с точкой
  {2mm}                       % расстояние между номером и текстом
  {}                           % перед заголовком можно добавить что-то
\titlespacing*{\section}{1.5cm}{*3}{*1}  % {отступ слева}{отступ сверху}{отступ снизу}

\captionsetup{justification=centering}
\renewcommand{\ttdefault}{cmtt}
\keepXColumns
\setlength{\parindent}{3em}
\newgeometry{left=2.5cm,right=2.5cm,top=3cm,bottom=2.5cm}

\renewcommand{\contentsname}{\hspace*{4em}{\smallСодержание}}
\setlength{\cftbeforesecskip}{0.5em}
\renewcommand{\cftsecleader}{\cftdotfill{\cftdotsep}}
\renewcommand{\cftsecaftersnum}{.\hspace{2mm}}

\title{Отчёт по культорологии}
\author{Выполнил студент группы М-РИТ-25-1, Станиславчук Сергей}

\begin{document}

\setstretch{1.5}
\noindent
\begin{minipage}[t]{0.01\textwidth}
    \includegraphics{~/programmer/++Programmer/5_1/LSTU.jpg}
\end{minipage}%
\begin{minipage}[c]{1.1\linewidth}
    \centering
    {\small
    МИНИСТЕРСТВО НАУКИ И ВЫСШЕГО ОБРАЗОВАНИЯ РФ \\[1em]
    ФЕДЕРАЛЬНОЕ ГОСУДАРСТВЕННОЕ БЮДЖЕТНОЕ \\
    ОБРАЗОВАТЕЛЬНОЕ УЧРЕЖДЕНИЕ ВЫСШЕГО ОБРАЗОВАНИЯ \\
    «ЛИПЕЦКИЙ ГОСУДАРСТВЕННЫЙ ТЕХНИЧЕСКИЙ УНИВЕРСИТЕТ»
    }
\end{minipage}%

\vspace{3em}

\noindent
\hspace{1em}%
\begin{minipage}[t]{0.8\textwidth}
\textbf{Институт} \hspace{2em} компьютерных наук \\
\textbf{Кафедра} \hspace{2em} автоматизированных систем управления
\end{minipage}


\vspace{6em}

\begin{center}
\textbf{ЗАДАНИЕ 6} \\
по дисциплине \\ <<Межкультурное взаимодействие>>
\end{center}

\vspace{6em}

\noindent
\begin{tabular*}{\textwidth}{@{\extracolsep{\fill}} l c r }
    Студент М-РИТ-25-1 & \makebox[4cm]{\hrulefill} & Станиславчук С.М. \\[2em]
Руководитель & \\
Ст. преподаватель & \makebox[4cm]{\hrulefill} & Болдырихин О.В. \\
\end{tabular*}

\vfill
\begin{center}
    Липецк, 2025
\end{center}

\thispagestyle{empty}

\newpage
\tableofcontents
\newpage

\section{Задание 1. Как воспринимается мир через призму культуры?}
\begin{sloppypar}
Культура задаёт человеку систему координат, через которую он интерпретирует
    происходящее. Она формирует представления о нормальном и ненормальном,
    хорошем и плохом, допустимом и запретном. Человек никогда не воспринимает
    реальность «как есть» — между ним и миром всегда стоит культурный фильтр.
    Через него мы распознаём значения, смыслы, символы, социальные роли и
    модели поведения.

    Кроме того, культура определяет способы реагирования на ситуации. То, что в
    одной культурной среде воспринимается как проявление уважения, в другой
    может читаться как слабость или грубость. Таким образом, культура не просто
    влияет на восприятие — она структурирует саму логику мышления, задаёт
    модели объяснения событий и определяет, какие аспекты реальности будут
    замечены, а какие проигнорированы.

    В результате человек видит мир не напрямую, а в интерпретации,
    сформированной воспитанием, традициями, нормами, коллективным опытом и
    социальными ожиданиями. Этот культурный контекст формирует индивидуальное и
    групповое мировосприятие, которое является устойчивым, но всё же поддаётся
    изменению при контакте с другими культурами.
\end{sloppypar}

\section{Задание 2. Какова роль языка в освоении действительности?}
\begin{sloppypar}
Язык — главный инструмент, через который человек усваивает и классифицирует
    реальность. С его помощью мы даём названия объектам и явлениям, определяем
    их свойства, выделяем существенные признаки и строим более сложные
    абстракции. Без языка невозможно формирование устойчивых понятий: он
    превращает разрозненные впечатления в структурированную систему знаний.

    Язык также формирует мышление. Его грамматика, словарный состав и логика
    определяют, какие идеи можно выразить легко, а какие требуют обходных
    формулировок. Каждая языковая система задаёт собственные «категориальные
    рамки», влияя на то, как человек различает время, пространство, действия,
    роли и отношения между предметами.

    Кроме того, язык обеспечивает передачу опыта между людьми и поколениями. С
    его помощью социальные группы формируют общую картину действительности,
    закрепляют нормы и правила, передают культурные смыслы. Роль языка
    — не просто отражение мира, а активное участие в его освоении и
    конструировании.
\end{sloppypar}

\section{Задание 3. Что означает понятие "картина мира"?}
\begin{sloppypar}
Картина мира — это целостная система представлений человека или общества о
    природе, обществе, человеке, ценностях, нормах и закономерностях бытия. Она
    объясняет устройство мира, связывает отдельные элементы реальности в
    логичную модель и задаёт ориентиры для поведения.

    В картину мира входят знания, убеждения, культурные традиции, идеология,
    религия, научные представления и личный опыт. Всё это объединяется в единое
    поле смыслов, благодаря которому человек понимает, что происходит вокруг, и
    способен предсказывать последствия своих действий.

    Также картина мира выполняет регулятивную функцию. Она подсказывает, как
    следует относиться к событиям, какие решения выбирать, как оценивать
    поступки свои и чужие. У разных обществ и эпох картины мира различаются:
    меняются научные взгляды, социальные нормы, ценности и способы объяснения
    необъяснимого.

    Картина мира — это внутренняя модель реальности, которая делает
    поведение осмысленным, а саму жизнь — структурированной и понятной.
\end{sloppypar}

\section{Задание 4. Что такое концептуальная картина мира?}
\begin{sloppypar}
Концептуальная картина мира представляет собой совокупность ключевых понятий,
    категорий и схем мышления, через которые человек или сообщество осмысливает
    реальность. Она отражает не только знания о предметах и явлениях, но и
    способы их интерпретации, связи между ними, причинно-следственные отношения
    и ценностные ориентиры.

    Эта картина формируется под влиянием культуры, образования, науки и личного
    опыта. В отличие от эмпирической картины мира, которая фиксирует
    наблюдаемые факты, концептуальная картина структурирует эти факты, создаёт
    абстрактные модели и позволяет строить прогнозы, делать выводы,
    систематизировать знания. Она задаёт рамки для понимания и объяснения
    явлений, определяя, какие вопросы будут считаться значимыми, а какие —
    второстепенными.

    Концептуальная картина мира выполняет когнитивную функцию: она помогает
    человеку ориентироваться в окружающей действительности, обрабатывать
    информацию и выстраивать целостное представление о мире, а также
    социально-культурную функцию, обеспечивая совместимость взглядов внутри
    общества.
\end{sloppypar}

\section{Задание 5. Что называется языковой картиной мира?}
\begin{sloppypar}
Языковая картина мира — это часть общей картины мира, отражённая через язык
    конкретного народа. Она включает в себя систему понятий, категорий,
    выражений, метафор и устойчивых словосочетаний, через которые носители
    языка осмысливают реальность. Язык не просто передаёт знания, он формирует
    мышление и определяет, как человек различает объекты, явления, действия и
    отношения между ними.

    В языковой картине мира отражены культурные ценности, исторический опыт и
    социальные нормы. Например, слова и выражения, обозначающие эмоции,
    природные явления или социальные роли, показывают, что именно в культуре
    считается значимым и как это интерпретируется. Она влияет на восприятие
    мира через способ категоризации и обозначения объектов, создавая уникальную
    «ленту» смыслов, присущую конкретному языковому сообществу.

    Языковая картина мира — это своеобразная когнитивная карта,
    обеспечивающая понимание действительности через призму языка, закрепляющая
    культурные представления и упорядочивающая опыт поколений.
\end{sloppypar}


\section{Задание 6. Что означает понятие "картина мира"?}
\begin{sloppypar}
Сложившаяся в сознании человека картина мира выступает фильтром восприятия,
    через который он интерпретирует информацию и оценивает события. Она
    определяет, какие явления будут замечены, а какие проигнорированы, какие
    события вызывают эмоциональный отклик, а какие остаются нейтральными.

    Картина мира задаёт критерии понимания и оценки: она формирует систему
    ценностей, ожиданий и норм, которые влияют на действия человека и его
    поведенческие реакции. Например, одно и то же событие в разных культурных
    или социальных контекстах будет восприниматься совершенно по-разному,
    поскольку интерпретация опирается на внутренние модели и убеждения.

    Кроме того, картина мира способствует предсказуемости поведения и упрощает
    обработку информации: человек опирается на привычные схемы, чтобы быстро
    ориентироваться в ситуации. Однако она также может ограничивать восприятие,
    создавая когнитивные стереотипы и препятствуя принятию новых знаний. В
    целом, картина мира формирует осмысленное восприятие действительности,
    направляет мышление и определяет индивидуальные и коллективные реакции на
    события.
\end{sloppypar}

\section{Задание 7. Насколько культура, в которой сформировалась личность
человека, оказывает влияние на его картину мира?} Культура является основным
фактором формирования картины мира человека. Она задаёт систему ценностей,
норм, традиций и моделей поведения, через которые индивид воспринимает
окружающую действительность. Через культурные практики и социальные институты
человек усваивает, что считается нормальным, значимым или допустимым в
обществе, а что — нет.

Влияние культуры проявляется и в способах интерпретации информации. Люди,
выросшие в разных культурных контекстах, по-разному понимают одни и те же
события, реагируют на них и оценивают их последствия. Таким образом, картина
мира не является универсальной для всех людей: она формируется под воздействием
конкретной культурной среды, влияя на восприятие, мышление и поведение
личности.

\section{Задание 8. Какую роль играет в этом процессе его родной язык?}
Родной язык служит основным инструментом усвоения и передачи культурного опыта.
Через язык человек получает знания о мире, усваивает социальные нормы и
культурные ценности, формирует категории восприятия. Лексика, грамматика и
устойчивые выражения языка задают определённые схемы мышления и способы
организации опыта.

Язык не только отражает мир, но и активно структурирует его в сознании
человека. Например, слова, обозначающие семейные роли, природные явления или
моральные категории, формируют представления о том, что важно, а что
второстепенно. Родной язык обеспечивает внутреннюю когнитивную карту, через
которую культура «проявляется» в индивидуальном восприятии и понимании
действительности.

\section{Задание 9. Почему время считается культурно-обусловленной категорией?}
Время не является исключительно объективной величиной; его восприятие и
интерпретация зависят от культуры. Разные общества различают прошлое, настоящее
и будущее по-разному, формируют уникальные ритуалы, праздники, циклы работы и
отдыха, что отражает их ценностные и мировоззренческие установки.

Культурная обусловленность времени проявляется в том, как люди планируют
действия, оценивают события и интерпретируют причинно-следственные связи.
Например, для одних культур характерна линейная концепция времени с
последовательным прогрессом, для других — циклическая, где события повторяются
и обретают смысл через ритуалы и традиции. Время становится не только
физическим, но и символическим измерением, через которое выражаются социальные
нормы и культурные ценности.

\section{Задание 10. Каким образом символика цвета позволяет увидеть картину мира, преломленную в сознании того или иного этноса, его глазами?}
\begin{sloppypar}
Цветовая символика — это элемент культурного кода, через который можно понять
    ценности, нормы и восприятие мира конкретного этноса. Разные культуры
    придают одному и тому же цвету различные значения: белый может
    символизировать чистоту и радость в одних обществах, а траур и смерть — в
    других.

    Символика цвета отражает эмоциональные и социальные установки, исторический
    опыт и религиозные представления народа. Анализ использования цветов в
    одежде, искусстве, ритуалах позволяет увидеть, какие категории важны для
    данного этноса, как он интерпретирует добро и зло, жизнь и смерть,
    священное и обыденное. Через цвет человек получает доступ к культурной
    логике, которая формирует восприятие мира и создает уникальную картину
    реальности для конкретной культурной группы.
\end{sloppypar}
\end{document}
